
% @Author: Vargas Hector <vargash1>
% @Date:   Saturday, June 25th 2016, 12:52:10 am
% @Email:  vargash1@wit.edu
% @Last modified by:   vargash1
% @Last modified time: Sunday, June 26th 2016, 10:22:53 pm
\documentclass{article}
% Package that renders tables nicely
\usepackage{longtable}
% Package that has high graphics rending capability
\usepackage{pgf}
\usepackage{amsmath}
\usepackage{amssymb}
\usepackage{siunitx}
\usepackage{graphicx}
\usepackage[version=3]{mhchem}
\usepackage[margin=1.5in]{geometry}
\pagenumbering{gobble}
\graphicspath { {./images/} }
\begin{document}
    %----------------------------------------------
    % Cover Page
    % A fancy cover page with centered text
    \thispagestyle{empty}
    \begin{center}
        \includegraphics[scale=0.3]{wit_logo}\\[1cm]
        Exam Writeup using  \LaTeX\\
        Professor G.Sirokman\\
        Hector Vargas\\
        CHEM 1100-4B
    \end{center}
    \pagebreak
    %----------------------------------------------

    %%%%%%%%%%%%%%%%%%%%%%%%%%%%%%%%%%%%%%%
    %%%%    ____                   __  %%%%
    %%%%   / __/_ _____ ___ _    <  /  %%%%
    %%%%  / _/ \ \ / _ `/  ' \   / /   %%%%
    %%%% /___//_\_\\_,_/_/_/_/  /_/    %%%%
    %%%%%%%%%%%%%%%%%%%%%%%%%%%%%%%%%%%%%%%

    %----------------------------------------------
    % Question 1 Exam 1
    \begin{center}
        \textbf{Exam 1 Part 1}\\
        \textit{Atomic Structure}
    \end{center}
    \textbf{1. There are two naturally occurring types of chlorine, ${}^{35}$ Cl $(34.969 \si{\atomicmassunit})$ and ${}^{37}$ Cl $(36.966 \si{\atomicmassunit})$}

    a) Given that the atomic weight of chlorine is $63.546 \si{\atomicmassunit}$, what are the abundances of ${}^{35}$Cl and ${}^{37}$Cl?

    b) What makes ${}^{35}$Cl different from ${}^{37}$Cl?

    Both of these isotopes of chlorine differ by the number of neutrons they both have. It's important to note that the number of protons must remain the same. A different number of protons results in a different element. The difference in neutrons is what contributes to the difference in weight.

    c)Which chlorine is regular chlorine and which one is the isotope?

    The term regular would be a bit incorrect. They are both regular chlorine, what makes them 'regular' is the fact that one of their isotopes is in higher abundance.

    \pagebreak
    %----------------------------------------------

    %----------------------------------------------
    % Question 2 Exam 1
    \begin{center}
        \textbf{Exam 1 Part 1}\\
        \textit{Unit Analysis}
    \end{center}
    \textbf{2. Consider the two different ways to arrange a large sheet of atoms}

    i)How many sulfur atoms(atomic radius $.180 \si{\nano\metre}$) can you fit on a square area $2.5 \si{\micro\metre} * 2.5 \si{\micro\metre}$ in arrangement(a)?

    ii)How many sulfur atoms(atomic radius $.180 \si{\nano\metre}$) can you fit on a square area $2.5 \si{\micro\metre} * 2.5 \si{\micro\metre}$ in arrangement(a)?

    iii)Which of these two layers is denser? By what factor?

    \pagebreak
    %----------------------------------------------

    %----------------------------------------------
    % Question 3 Exam 1
    \begin{center}
        \textbf{Exam 1 Part 1}\\
        \textit{Atoms}
    \end{center}
    \textbf{3. Consider the modern model of an atom}

    a) Draw a diagram of the structure of the atom. Indicate what type of particles exist in each region of the atom.

    b)Describe the properties of each of the three particles in the atom.

    c) What is the approximate diameter of an atom?
    \pagebreak
    %----------------------------------------------

    %----------------------------------------------
    % Question 4 Exam 1
    %
    \begin{center}
        \textbf{Exam 1 Part 2}\\
        \textit{Acid/Base}
    \end{center}
    \textbf{4. Consider the following questions about acids and bases.}

    a) Note for each compound if it is\\
    \begin{enumerate}
        \item an acid or a base
        \item weak or strong
    \end{enumerate}

    \ce{NH_{3}} \qquad \ce{HCl}

    \ce{KI} \qquad \ce{Ca(OH)_{2}}

    b) Predict the product of the following reaction. Make sure the reaction is balanced.

    c) Write a net ionic equation for b).

    d) Identify the acid and the base in the following reaction(remember the Bronsted Lowry definition)
    $$\ce{CH_{3}COOH} + \ce{NH_{3}} \rightarrow \ce{CH_{3}COO^{-}} + \ce{NH_{4}^{+}}$$
    \pagebreak

    %----------------------------------------------

    %----------------------------------------------
    % Question 5 Exam 1
    %
    \begin{center}
        \textbf{Exam 1 Part 2}\\
        \textit{Reduction/Oxidation}
    \end{center}
    \textbf{5. In the following reactions}
    \begin{enumerate}
        \item Identify the oxidation state of the elements involved
        \item Indicate what element is being oxidized and what is being reduced
        \item Indicate how many electrons were transferred in the reaction
    \end{enumerate}

    a) $\ce{2 C_{2}H_{2}} + \ce{3 O_{2}} \rightarrow \ce{2CO_{2}} + \ce{2 H_{2}O}$

    b) $\ce{Cu^{2+}} + \ce{Zn} \rightarrow  \ce{Zn^{2+}} + \ce{Cu}$

    c) $\ce{P_{4}} + \ce{10 HClO} + \ce{6 H_{2}O} \rightarrow \ce{4 H_{3}PO_{4}} + \ce{10 HCl}$
    \pagebreak
    %----------------------------------------------

    %----------------------------------------------
    % Question 6 Exam 1
    %
    \begin{center}
        \textbf{Exam 1 Part 2}\\
        \textit{Galvanic Series}
    \end{center}
    \textbf{6. You are handed three metal samples. You are told that the three samples are silver, cobalt, and manganese, but that it isn't known which is which. You have available to you some zinc nitrate solution, nickel nitrate solution, and concentrated hydrochloric acid}

    Describe in detail, with specifics, how you could use the chemicals available to identify the three metal samples. Feel free to make use of the galvanic series below.

    \pagebreak
    %----------------------------------------------

    %----------------------------------------------
    % Question 7 Exam 1
    %
    \begin{center}
        \textbf{Exam 1 Part 3}\\
        \textit{Stoichiometry w/ limiting reagent}
    \end{center}
    \textbf{7. One process used in gold mines to recover gold from mined rock is as shown below}
    $$  \ce{Au} +  \ce{NaCN} +   \ce{O_{2}} +  \ce{H_{2}O} \rightarrow  \ce{NaAu(CN)_{2}} +  \ce{NaOH}$$

    a) Balance the reaction

    b) A small test batch is run in $10 \si{\milli\liter}$($10 \si{\milli\liter}$ is $10 \si{\gram}$ of water). The sample in the reaction contains $0.0178 \si{\gram}$ of gold. It is treated with $0.0115 \si{\gram}$ of $\ce{NaCN}$. $\ce{O_{2}}$ is present in excess. What is the limiting reagent in this test reaction?

    c) What mass of $\ce{NaAu(CN)_{2}}$ would be produced by the reaction described by b?

    \pagebreak
    %----------------------------------------------

    %----------------------------------------------
    % Question 8 Exam 1
    %
    \begin{center}
        \textbf{Exam 1 Part 3}\\
        \textit{Stoichiometry With Story}
    \end{center}
    \textbf{8. Back in the day, aluminum was a highly precious metal. It was first isolated in the early 1800s. It's synthesis involved the reaction of aluminum chloride with potassium}
    $$ \ce{3K} + \ce{AlCl_{3}} \rightarrow \ce{Al} + \ce{3KCl}$$

    The Baron Von Markov has decided that he wants a full aluminum place setting, and you have been pout in charge of making the required aluminum. You have to produce $520 \si{\gram}$ of aluminum using the reaction shown above, which work in $87\%$ yield. How much aluminum chloride($\ce{AlCl_{3}}$) and potassium($\ce{K}$) will you need to be able to produce the needed amount of aluminum?

    \pagebreak
    %----------------------------------------------

    %----------------------------------------------
    % Question 9 Exam 1
    %
    \begin{center}
        \textbf{Exam 1 Part 3}\\
        \textit{Stoichiometry With Solutions Chemistry}
    \end{center}
    \textbf{9. Early lighter than air balloons often used hydrogen gas for lift. This hydrogen gas was generally produced from the reaction of iron and sulfuric acid as shown here:}
    $$\ce{Fe} + \ce{H_{2}SO_{4}} \rightarrow \ce{FeSO_{4}} + \ce{H_{2}}$$

    a) Is this reaction an acid base reaction or an oxidation reduction reaction(or neither?). Justify your answer.

    b) In a test reaction, reacting $100 \si{\gram}$ of iron with excess sulfuric acid produced $1.45 \si{\mole}$ of hydrogen gas. What is the percent yield of this reaction?

    c) A balloon with the volume of $100 \si{\liter}$ needs to be filled with hydrogen. How much iron and sulfuric acid will be needed to produce the required amount of hydrogen gas to fill the $100 \si{\liter}$ balloon? To help, 1 $\si{\mole}$ of gas fills $24 \si{\liter}$ of volume. \textbf{Make sure to account for the yield calculated in b!}

    \pagebreak
    %----------------------------------------------

    %----------------------------------------------
    % Bonus Exam 1
    %
    \begin{center}
        \textbf{Exam 1 Bonus}\\
    \end{center}

    A) Describe the phenomenon due to which the region of the atmosphere we call the sky is blue at noon on a sunny day. (aka Why is the sky blue?)

    B) We have talked about two ways to categorize compounds so far, molecules, and ionic compounds. How are these two types of compounds different from each other? How does difference influence the way they behave in solution?

    \textbf{Answer}\\
    These two different types of compounds primarily differ in the method in which they bond. Ionic compounds transfer electrons, whereas molecules share electrons.

    \pagebreak
    %----------------------------------------------

    %%%%%%%%%%%%%%%%%%%%%%%%%%%%%%%%%%%%%%%
    %%%%    ____                  ____ %%%%
    %%%%   / __/_ _____ ___ _    |_  | %%%%
    %%%%  / _/ \ \ / _ `/  ' \  / __/  %%%%
    %%%% /___//_\_\\_,_/_/_/_/ /____/  %%%%
    %%%%%%%%%%%%%%%%%%%%%%%%%%%%%%%%%%%%%%%

    %----------------------------------------------
    % Question 1 Exam 2
    %
    \begin{center}
        \textbf{Exam 2 Part 1}\\
        \textit{Stoichiometry}
    \end{center}
    \textbf{1. The spores of the hive mind have just landed on a planet, and you have begun your hungry search for more raw matter to build a new hive cluster. OF special interest to you are local minerals, which are predominantly unobtanium trioxide ($\ce{UbO_{3}}$). The stars in the Koprulu sector are capable of making super heavy atoms, so unobtanium($\ce{Ub}$) has an atomic weight of $558 \dfrac{\si{\gram}}{\si{\mol}}$}

    You need to hatch some Zerglings fast to kill some of the more violent local fauna, like the human farmers over the next hill. A zergling egg requires $400 \si{\kilo\gram}$ of pure unobtanium to produce, and each egg hatches two zerglings. Carbon(readily available) and unobtanium trioxide react to form unobtanium and carbon dioxide. This reaction has a yield of $85\%$ for unobtanium.

    How many kilograms of minerals and how many kilograms of carbon do you need to make 10 zerglings?

    \pagebreak
    %----------------------------------------------

    %----------------------------------------------
    % Question 2 Exam 2
    %
    \begin{center}
        \textbf{Exam 2 Part 1}\\
        \textit{Unit Analysis}
    \end{center}
    \textbf{2. Consider the two different ways to arrange a large sheet of atoms}

    i) How many gold atoms(atomic radius $.144 \si{\nano\metre}$) can you fit on a square area $1.5 \si{\micro\metre} * 1.5 \si{\micro\metre}$ in arrangement(a)?

    ii) How many gold atoms(atomic radius $.144 \si{\nano\metre}$) can you fit on a square area $1.5 \si{\micro\metre} * 1.5 \si{\micro\metre}$ in arrangement(a)?

    iii) Calculate the density $(\dfrac{\text{mass}}{\text{volume}})$ of a single atom layer of arrangement a)

    \pagebreak
    %----------------------------------------------

    %----------------------------------------------
    % Question 3 Exam 2
    %
    \begin{center}
        \textbf{Exam 2 Part 1}\\
        \textit{Red/Ox}
    \end{center}
    \textbf{3. Jewelers have numerous tests to be able to tell whether a piece of jewelry is fake. For instance, lead is often used in making fake gold items. A jeweler is handed a gold ring, but is suspicious as to whether it is gold or not. HE uses a tool to make a tiny scratch on the surface of the ring, and drips some concentrated hydrochloric acid($\ce{H^{+}}$) on it. He then observes that bubbles are formed in the acid solution. Use the table below to explain the chemistry that was observed if the ring was made of lead.}

    Although gold is not damaged by hydrochloric acid, it can be corroded by aqua regia, a mixture of nitric acid and hydrochloric acid. Below is the reaction between gold and aqua regia. Confirm that it is a red/ox reaction. Identify what element is oxidized and what is reduced. Indicate how many electrons are transferred.
    $$\ce{Au} + \ce{HNO_{3}} + \ce{4 HCl} \rightarrow \ce{HAuCl_{4}} + \ce{2 H_{2}O} + \ce{NO}$$

    \pagebreak
    %----------------------------------------------

    %----------------------------------------------
    % Question 4 Exam 2
    %
    \begin{center}
        \textbf{Exam 2 Part 2}\\
        \textit{Photon Energies}
    \end{center}
    \textbf{4. The Imperial Star Destroyer, His Left Hand, has departed space dock and is cruising around Tatooine. It spots the Rebel Frigatte, Valorous, and opens fire. His Left Hand's turbo laser batteries fires lasers with photons of wavelength $551 \si{\nano\meter}$.}

    a) If \textit{Valorous} is hit by a $10 \si{\mega\joule}$ blast of energy from the laster fire, how many photons hit the Valorous?

    $$ h = 6.626 * 10^{-34} \si{\joule} \cdot \si{\second}$$
    $$ c = 3.00 * 10^{8} \frac{\si{\meter}}{\si{\second}} $$

    \textbf{Answer}\\
    $$ E_{\text{photon}} = \dfrac{hc}{5.51 * 10^{-7} \si{\meter}} = 3.8900196 * 10^{-19} \si{\joule} $$

    b) \textit{His Left Hand} is using a copper based laser emitter. What is happening inside the copper atoms to emit photons?

    \pagebreak
    %----------------------------------------------

    %----------------------------------------------
    % Question 5 Exam 2
    %
    \begin{center}
        \textbf{Exam 2 Part 2}\\
        \textit{Electron Configuration}
    \end{center}
    \textbf{5. Consider the following questions relating to electron configurations of atoms.}

    a) Write the electron configurations of the following elements:
    \begin{enumerate}
        \item S\\
        \textbf{Answer: }
        $[\ce{Ne}]3s^{2}3p^{4}$
        \item N\\
        \textbf{Answer: }
        $[\ce{He}]2s^{2}2p^{3}$
    \end{enumerate}

    b) Draw an energy diagram for the electron configuration of oxygen (with lines for each orbital and electrons as arrows). What ion is the most likely ion that oxygen will form? Explain based on your electron configuration.

    c) Using quantum numbers, explain why the $n = 2$ shell has an $s$ and $p$ sub shell but no $d$ or $f$ sub shells.
    \pagebreak
    %----------------------------------------------

    %----------------------------------------------
    % Question 6 Exam 2
    %
    \begin{center}
        \textbf{Exam 2 Part 2}\\
        \textit{Orbitals}
    \end{center}
    \textbf{6. Consider the following questions relating to electron configurations of atoms.}

    a)Draw a digram, showing the probability of finding the electron on the $y$ axis, and distance from the nucleus on the $x$ axis for a $2s$ orbital

    b)Draw a diagram, showing the probability of finding the electron on the $y$ axis, and distance from the nucleus on the $x$ axis for a $2p$ orbital.

    c)Draw a picture of a $2p$ orbital

    d)In your diagrams and pictures above. Indicate the locations of any nodes. How many nodes does each type of orbital have?

    e)What appears to be a common theme between $2s$ and $2p$ orbitals based on your answers in d?


    \pagebreak
    %----------------------------------------------

    %----------------------------------------------
    % Question 7 Exam 2
    %
    \begin{center}
        \textbf{Exam 2 Part 3}\\
        \textit{Atomic Size}
    \end{center}
    \textbf{7. Consider the following questions relating to atomic size and lattice energy.}

    a) Put the following sets of atoms in order of increasing size and \textbf{explain} your choice for ordering
    $$\ce{O}, \ce{N}, \ce{F}, \ce{C}$$

    b) There are two components to changes in lattice energies. The values below are an example of each. What is suggested by each trend? What changes appear to have the greater influence on lattice energy? Explain in detail.
    \begin{enumerate}
        \item $\ce{LiF}$ $1060 \dfrac{\si{\kilo\joule}}{\si{\mole}} \qquad$ $\ce{CsI}$  $600 \dfrac{\si{\kilo\joule}}{\si{\mole}}$
        \item  $\ce{LiF}$ $1060 \dfrac{\si{\kilo\joule}}{\si{\mole}} \qquad$ $\ce{MgO}$  $3900 \dfrac{\si{\kilo\joule}}{\si{\mole}}$
    \end{enumerate}
    \textbf{Answer}\\

    1. $\ce{LiF}$ has a greater lattice energy due to the size difference. The increased atomic radii of $\ce{Cs}$ and $\ce{I}$ results in less lattice energy. Lithium and Fluorine are smaller relative to Cesium and Iodine, resulting in a tighter pull from the nucleus to the electrons. Thus, $\ce{LiF}_{\text{Lattice Energy}} > \ce{CsI}_{\text{Lattice Energy}}$

    2. In contrary to the previous question. The charge of each element has more influence on the lattice energy. $\ce{Mg^{2+}}$ and $\ce{O^{2-}}$ have a greater net charge difference than $\ce{Li^{+}}$ and $\ce{F^{-}}$
    \pagebreak
    %----------------------------------------------

    %----------------------------------------------
    % Question 8 Exam 2
    %
    \begin{center}
        \textbf{Exam 2 Part 3}\\
        \textit{Bond Enthalpy}
    \end{center}
    \textbf{8. Currently the vast majority of commercially used hydrogen is generated according to the following reaction.}
    $$\ce{CH4} + \ce{H2O} \rightarrow \ce{CO2} + \ce{H2}$$
    \begin{center}
        $\ce{C-H}$ $413 \frac{\si{\kilo\joule}}{\si{\mol}} \quad$
        $\ce{C-O}$ $358 \frac{\si{\kilo\joule}}{\si{\mol}} \quad$
        $\ce{O-O}$ $146 \frac{\si{\kilo\joule}}{\si{\mol}} \quad$
        $\ce{C=O}$ $799 \frac{\si{\kilo\joule}}{\si{\mol}}$\\[.2cm]
        $\ce{O-H}$ $463 \frac{\si{\kilo\joule}}{\si{\mol}} \quad$
        $\ce{C-C}$ $348 \frac{\si{\kilo\joule}}{\si{\mol}} \quad$
        $\ce{O=O}$ $495 \frac{\si{\kilo\joule}}{\si{\mol}} \quad$
        $\ce{H-H}$ $436 \frac{\si{\kilo\joule}}{\si{\mol}}$
    \end{center}

    a) What is the enthalpy change($\Delta H$) for this reaction per $\ce{CO2}$? Is energy released or absorbed?

    b) There are a variety of means of recovering energy from $\ce{H2}$, either from giving heat through combustion or electricity electrochemically from a fuel cell. What is the energy released in this reaction?
    $$\ce{H2} + \ce{O2} \rightarrow \ce{2H2O}$$

    c) Combining a) and b), do some stoichiometry to find the net energy released per $\ce{CO2}$ produced.
    \pagebreak
    %----------------------------------------------

    %----------------------------------------------
    % Question 9 Exam 2
    %
    \begin{center}
        \textbf{Exam 2 Part 3}\\
        \textit{Bond Polarity}
    \end{center}
    \textbf{8. Consider the following questions about bonding and polarity.}

    a) Order the following bonds from most to least polar. Show your calculations.
    \begin{enumerate}
        \item $\ce{F-F}$ in $\ce{F2}$
        \item $\ce{S=C}$ in $\ce{CS2}$
        \item $\ce{O-H}$ in $\ce{H2O}$
        \item $\ce{Cl-H}$ in $\ce{HCl}$
        \item $\ce{B-F}$ in $\ce{BF3}$
        \item $\ce{Na-Cl}$ in $\ce{NaCl}$
    \end{enumerate}

    b)  Categorize each of the bonds in a) as polar covalent, non polar covalent, or ionic.

    c) What does it mean for a bond to be polar? How is this different from a bond being ionic?
    \pagebreak
    %----------------------------------------------

    %----------------------------------------------
    % Bonus Exam 2
    %
    \begin{center}
        \textbf{Exam 2 Bonus}\\
    \end{center}

    A) Describe the phenomenon due to which the region of the atmosphere we call the sky is blue at noon on a sunny day. (aka Why is the sky blue?)

    B) Use Molecular Orbitals to explain whether or not $\ce{He2+}$ molecule should be able to exist.

    \textbf{Answer}\\
    $$\text{Bond Order} = \dfrac{2 - 0}{2} = 1$$
    It would be reasonable for $\ce{He2+}$ to exist. A bond order of 0 would suggest it's too unstable.

    \pagebreak
    %----------------------------------------------

    %%%%%%%%%%%%%%%%%%%%%%%%%%%%%%%%%%%%%%%
    %%%%    ____                  ____ %%%%
    %%%%   / __/_ _____ ___ _    |_  / %%%%
    %%%%  / _/ \ \ / _ `/  ' \  _/_ <  %%%%
    %%%% /___//_\_\\_,_/_/_/_/ /____/  %%%%
    %%%%%%%%%%%%%%%%%%%%%%%%%%%%%%%%%%%%%%%

    %----------------------------------------------
    % Question 1 Exam 3
    %
    \begin{center}
        \textbf{Exam 3 Part 1}\\
        \textit{Stoichiometry}
    \end{center}
    \textbf{1. When hydrogen sulfide gas ($\ce{H2S}$) is bubbled into a solution of sodium hydroxide ($\ce{NaOH}$), the reaction forms sodium sulfide ($\ce{Na2S}$) and water. How many grams of sodium sulfide are formed if $3.50 \si{\gram}$ of hydrogen sulfide is bubbled into a solutions containing $4.20 \si{\gram}$ of sodium hydroxide, assuming that the sodium sulfide is made in $73\%$ yield?.}

    \pagebreak
    %----------------------------------------------

    %----------------------------------------------
    % Question 2 Exam 3
    %
    \begin{center}
        \textbf{Exam 3 Part 1}\\
        \textit{Spectral Lines}
    \end{center}
    \textbf{2. The Rydberg equation can be used to predict the spectral emissions of hydrogen atoms. $R_{H} = 1.097 * 10^{7}  \si{\meter^{-1}}$; $c = 3.00 * 10^{8} \si[per-mode=symbol]{\meter\per\second}$; $h = 6.626 * 10^{-34} \si{\joule} \si{\second}$.}

    \textbf{a)  The Balmer series is the series of emissions arising from the fall of an electron from a higher orbital to the $n = 2$ orbital. Calculate the wavelength of the lowest energy transition in the Balmer series.}\\
    \textbf{\textit{Solution}}
    $$\dfrac{1}{\lambda} = R_{H}(\dfrac{1}{n^{2}_{f}} - \dfrac{1}{n^{2}_{i}}) \rightarrow \dfrac{1}{\lambda} = R_{H}(\dfrac{1}{2^{2}} - \dfrac{1}{3^{2}})$$
    $$\downarrow$$
    $$\dfrac{1}{\lambda} = 1523611.11^{-1} \rightarrow \lambda = 656 \si{\nano\metre}$$
    It's important to note that this wavelength is well within our viewing range!

    \textbf{b) The Lyman series is a set of emissions arising from the fall of an electron from a higher orbital to the $n = 1$ orbital. Do you expect emissions in this series to be higher in energy or lower than the Balmer series? Explain.}\\
    \textbf{\textit{Solution}}
    $$\dfrac{1}{\lambda} = R_{H}(\dfrac{1}{1^{2}} - \dfrac{1}{2^{2}}) \rightarrow 1.215 * 10^{-7} = 121 \si{\nano\meter}$$
    $$E = \dfrac{hc}{\lambda} $$
    $$ \downarrow $$
    $$E_{656 \si{\nano\meter}} = \dfrac{hc}{6.56 * 10^{-7}}  = 3.030 * 10^{-19} \si{\joule}$$
    $$E_{121 \si{\nano\meter}} = \dfrac{hc}{1.215 * 10^{-7}} = 1.636 * 10^{-18} \si{\joule}$$

    Shorter wavelengths have higher energy. Which is why gamma rays are dangerous, but transitions from higher orbital to lowest increase in energy. Such as $n=5$ to $n=2$. Thus, these emissions would be higher. The real contributor to more energy is higher $n$ orbital drop.

    \pagebreak
    %----------------------------------------------

    %----------------------------------------------
    % Question 3 Exam 3
    %
    \begin{center}
        \textbf{Exam 3 Part 1}\\
        \textit{Bond Enthalpy}
    \end{center}
    \textbf{4. Ethanol is a common additive to gasoline nowadays. Ethanol, however, is said to be less energy dense.\\ Here are a few energy bonds:}
    \begin{center}
        $\ce{C-H}$ $413 \  \si[per-mode=symbol]{\kilo\joule\per\mol} \quad$
        $\ce{C-O}$ $358 \ \si[per-mode=symbol]{\kilo\joule\per\mol} \quad$
        $\ce{O-O}$ $146 \  \si[per-mode=symbol]{\kilo\joule\per\mol} \quad$
        $\ce{C=O}$ $799 \  \si[per-mode=symbol]{\kilo\joule\per\mol}$\\[.2cm]
        $\ce{O-H}$ $463 \  \si[per-mode=symbol]{\kilo\joule\per\mol} \quad$
        $\ce{C-C}$ $348 \  \si[per-mode=symbol]{\kilo\joule\per\mol} \quad$
        $\ce{O=O}$ $495 \  \si[per-mode=symbol]{\kilo\joule\per\mol}$
    \end{center}

    \textbf{a) How much energy is released per mole of ethanol($\ce{C2H5OH}$)?}

    \textbf{b) How much energy is released per liter of ethanol burnt?(density = $.789 \si[per-mode=symbol]{\gram\per\milli\liter}$)?}
    \pagebreak
    %----------------------------------------------

    %----------------------------------------------
    % Question 4 Exam 3
    %
    \begin{center}
        \textbf{Exam 3 Part 2}\\
        \textit{MO Theory}
    \end{center}
    \textbf{4.}

    \textbf{a) Draw a molecular orbital digram for the molecule $\ce{F2}$}

    \textbf{b) What is the Bond Order of $\ce{F2}$(show calculations)? Does this match what is predicted by Lewis structures(You will have to draw a Lewis structure to convince me)?}
    $$ \dfrac{10 - 8}{2} = 1$$
    $$ \ce{F-F}$$

    \textbf{c) Is $\ce{F2}$ diamagnetic or paramagnetic? How will it react to a magnetic field?}

    $\ce{F2}$ is paramagnetic due to a set of  unpaired electrons, it will be attracted to a magnetic field.
    \pagebreak
    %----------------------------------------------

    %----------------------------------------------
    % Question 5 Exam 3
    %
    \begin{center}
        \textbf{Exam 3 Part 2}\\
        \textit{Heating Curves}
    \end{center}
    \textbf{5. \textit{Valorous}, the rebel frigate is having a luck day. \textit{His Left Hand's} (Imperical Cruiser) shields are down, so a $15 \si{\giga\joule}$ blast of energy hits the side of the ship. Let us assume that this ship was made with low grade materials like steel(we'll approximate it's characteristics with iron). Calculate what the maximum amount of steel/iron this shot could vaporize is. Assume the side of the ship is at $300 \si{\kelvin}$ due to the heat generated inside the Cruiser. The melting point of iron is $1811 \si{\kelvin}$, the boiling point is $3134 \si{\kelvin}$, the specific heat of solid iron is $.445 \si[per-mode=symbol]{\joule\per\gram\kelvin}$. The specific heat of liquid iron is $.611 \si[per-mode=symbol]{\joule\per\gram\kelvin}$, the heat of fusion for iron is $13.8 \si[per-mode=symbol]{\kilo\joule\per\mole}$ and the heat of vaporization for iron is $340 \si[per-mode=symbol]{\kilo\joule\per\mole}$. Build a strategy to solve the problem. As one option, try to find the energy to boil a small amount of the material and go from there.}
    \pagebreak
    %----------------------------------------------

    %----------------------------------------------
    % Question 6 Exam 3
    %
    \begin{center}
        \textbf{Exam 3 Part 2}\\
        \textit{Molecular Orbitals}
    \end{center}
    \textbf{6. Below is a presented the molecular orbital diagram for $\ce{Li2}$}

    \textbf{a) What is the bond order for this molecule?}

    \textbf{\textit{Solution}}
    $$\text{Bond Order} = \dfrac{4 - 2}{2} = 1$$

    \textbf{b) Which of the molecular orbitals will have nodal planes? What does this imply about these orbitals(the ones with nodal planes?)}

    \textbf{c) The $\ce{Li2}$ molecule is hit by a photon, and one of the $\sigma_{2s}$ electrons is promoted to the $\sigma_{2s}^{*}$ orbital. Would this be a stable molecule? }

    \textbf{\textit{Solution}}
    $$\text{New Bond Order} = \dfrac{2 - 2}{2} = 0$$
    The bond order is now 0, this leads to the molecule being unstable and rather unlikely to exist.


    \textbf{d) The strength of a $\ce{Li-Li}$ bond in $\ce{Li2}$ is $215 \si[per-mode=symbol]{\kilo\joule\per\mole}$. What is the wavelength of the photon in question c?}
    $$ c = 3.0 * 10^{8} \si[per-mode=symbol]{\meter\per\second} \quad h = 6.626 * 10^{-34} \si{\joule\second}$$

    \pagebreak
    %----------------------------------------------

    %----------------------------------------------
    % Question 7 Exam 3
    %
    \begin{center}
        \textbf{Exam 3 Part 3}\\
        \textit{Acids/Bases}
    \end{center}
    \textbf{7.}

    \textbf{a) What is the pH of a $.0034 M$ solution of         $\ce{NaOH}$?}

    \textbf{b) What is the pH of a $.0034 M$ solution of         $\ce{CH3COOH}$? ($k_{a} = 1.8 * 10^{-5}$)}
    $$\ce{CH3COOH <=> CH3COO- + H+}$$

    \textbf{c) what is the pH of a solution made from equal volumes of the solutions in a and b?(Consider what you are starting off with after the two solutions react. Same as $k_{a}$ as above applies.)}
    $$\ce{CH3COOH + OH- <=> CH3COO- + H2O}$$

    \pagebreak
    %----------------------------------------------

    %----------------------------------------------
    % Question 8 Exam 3
    %
    \begin{center}
        \textbf{Exam 3 Part 3}\\
        \textit{Intermolecvular Forces}
    \end{center}
    \textbf{8. Consider the following questions about Van Der Waals forces.(Hint Draw lewis structures)}

    \textbf{a) Ethane $(\ce{C2H6})$ and formaldehyde ($\ce{CH2O}$) have the same molar mass and similar size. Formaldehyde, however, boils at $-19\,^{\circ}\mathrm{C}$, and ethane boils at $-89\,^{\circ}\mathrm{C}$. Explain this difference in terms of intermolecular forces.}

    The difference in boiling point suggests that formaldehyde has a particular type of intermolecular force that ethane does not have. Ethane's lewis structure is simply 2 carbons surrounded by hydrogen. Thus, it has no dipole-dipole interaction. We know that all compounds have London dispersion forces, thus it leaves us with Formaldehyde either having dipole-dipole attractions or h-bonds. However, formaldehyde cannot have h-bonds as it does not meet the criteria for doing so. Thus we know that formaldehyde has dipole-dipole intermolecular forces while ethane does not.  We can confirm this with the usage of Lewis structures. As you can see, the oxygen's charge is the primary reason for the dipole-dipole charge.

    \textbf{b) Methane $(\ce{CH4})$ and water ($\ce{H2O}$) have the same molar mass and similar size. None the less, methane boils at $-161\,^{\circ}\mathrm{C}$  whereas water boils at $100\,^{\circ}\mathrm{C}$. Explain this  dramatic difference in terms of intermolecular forces.}

    Unlike Methane, water has a hydrogen bond. These types of bonds are much stronger than London dispersion forces and dipole-dipole forces. Methane only has London dispersion forces, similar to Ethane in the previous example. This would explain why $\ce{HF}$ is known for having such a strong bond.
    \pagebreak
    %----------------------------------------------

    %----------------------------------------------
    % Question 9 Exam 3
    %
    \begin{center}
        \textbf{Exam 3 Part 3}\\
        \textit{Equilibrium}
    \end{center}
    \textbf{9. Consider the reactions:}
    $$ \ce{2NO} (\si{\gram}) \ce{<=>} \ce{N2}(\si{\gram}) + \ce{O2}(\si{\gram})$$
    $$\text{The K of the reaction is } 2.4*10^{3} \text{ at }  2000\,^{\circ}\mathrm{C}$$

    If $0.256 \si{\mole}$ of $\ce{NO}$ is introduced into a $1.35 \si{\liter}$ reactor and is allowed to come to equilibrium what will be the final concentrations of $\ce{NO }$, $\ce{N2}$, and $\ce{O2}$ be?

    \pagebreak
    %----------------------------------------------

    %----------------------------------------------
    % Bonus Exam 3
    %
    \begin{center}
        \textbf{Exam 3 Bonus}\\
    \end{center}

    A) Describe the phenomenon due to which the region of the atmosphere we call the sky is blue at noon on a sunny day. (aka Why is the sky blue?)

    B) The synthesis of biodiesel takes advantage of Le Chatelier's principle in two ways. Describe Le Chatelier's principle and how it helps generate higher purity biodiesel in the lab.

    %----------------------------------------------

\end{document}
